% ===================================================================
%  TABLO No 11 — Şehir ve yapıları. — Yangın.
%  Traduction 2026 d'apres texto/io, controlee sur texto/fr et
%  texto/hi. Consigne : voir texto/tr/10-tablo-01.tex.
% ===================================================================

%%K t11-apar-1 apar
\VUsaut{2.0mm}
\VUcentre{13.2pt}{90}{ÜÇÜNCÜ DİZİ}
\VUsaut{3.0mm}
\VUfilet{16mm}
\VUsaut{5.6mm}
\VUtitre{15.0pt}{60}{TABLO No 11}
\VUsaut{1.4mm}
\VUpk{11.4pt}{40}{Şehir ve yapıları. --- Yangın.}
\VUsaut{2.2mm}

%%K t11-tit-1 sub
\VUpk{10.2pt}{40}{Dış mahalle.}
\VUsaut{3.0mm}

%%K t11-c1-01-1 p
1. --- Okyanustan yüz kilometre uzakta büyük bir şehirde otururum, ki yine de
birinci sınıf bir \VUgras{limandır} \textsuperscript{(1)}. Yanından akan ırmak
dört yüz metre geniştir ve pek güzel bir demirleme yeri yapar.

%%K t11-c1-02-1 p
2. --- \VUgras{Rıhtımlar} \textsuperscript{(2)} boyunca uzanan bütün bölüm tam
denizci ve sanayici bir görünüş taşır, ve içinde yalnız denizcilerle işçilerin
oturduğu bir \VUgras{dış mahalle} \textsuperscript{(3)} yapar. Bunlar
\VUgras{fabrikalarda} \textsuperscript{(4)} ve \VUgras{havagazı fabrikasında}
\textsuperscript{(5)} çalışır, ki yüksek \VUgras{bacası} \textsuperscript{(6)}
ile \VUgras{gaz deposu} uzaktan görünür.

%%K t11-c1-03-1 p
3. --- Orta çağda şehir surluydu, ve surlarından birkaç kalıntı bugün de
bulunur, özellikle o \VUgras{kapı} \textsuperscript{(8)}, ki eski \VUgras{iç
kalenindir} \textsuperscript{(7)}, ve iki yanında iki \VUgras{kulesi}
\textsuperscript{(9)} vardır; onlar da şimdi \VUgras{hapishane} olarak
kullanılır.

%%K t11-c1-04-1 p
4. --- Pek eski görünen bütün bu \VUgras{mahallede} yalnız bir yapı görece
yenidir: \VUgras{gümrük binası} \textsuperscript{(10)}. Rıhtımla o küçük
\VUgras{sokağın} \textsuperscript{(11)} arasında durur, ki eski
\VUgras{belediye konağına} \textsuperscript{(12)} gider. Bu son yapı, eski
şehrin üstüne çöken o kocaman \VUgras{çan kulesinden} \textsuperscript{(13)}
kolayca tanınır.

%%K t11-c1-05-1 p
5. --- Sokağın öte yanında gümrük binasının biçemindeki bir yapı
\VUgras{yükselir}: \VUgras{ticaret odasıdır} \textsuperscript{(14)}. Bir yüzü
küçük bir alana açılır, ki orada \VUgras{valilik} \textsuperscript{(15)}
bulunur. Şehrimiz, gerçekten, başkent olmasa da önemli bir il merkezidir.

%%K t11-tit-2 sub
\VUsaut{5.0mm}
\VUpk{10.2pt}{40}{Şehrin yukarı bölümü.}
\VUsaut{3.4mm}

%%K t11-c1-06-1 p
6. --- Epeyce büyük olan garnizon, geniş ve pek sağlıklı \VUgras{kışlalarda}
\textsuperscript{(16)} oturur; onlar \VUgras{şehir parkının}
\textsuperscript{(17)} parmaklığına dek uzanır. Bu parka giden
\VUgras{geniş cadde} \textsuperscript{(19)}, \VUgras{tasarruf sandığının}
\textsuperscript{(18)} arkasından geçip \VUgras{katedralin}
\textsuperscript{(20)} alanına çıkar.

%%K t11-c1-07-1 p
7. --- Bütün bu yapılar küçük bir tepede basamak basamak yükselir; orada büyük
\VUgras{lisemiz} \textsuperscript{(21)} de bulunur.

%%K t11-c1-07-2 p
Büyük caddeyle birleşen o hafif eğimli yoldan inerseniz, \VUgras{adliyeye}
\textsuperscript{(22)} varırsınız, ki önünde hoş bir \VUgras{kamu bahçesi}
\textsuperscript{(26)} vardır. Bu \VUgras{alan} pek büyük değildir, ama şehrin
ortasında bir yeşillik ve serinlik vahası yapar. Bunun için şehirliler oraya
bol bol gelir; \VUgras{kahvenin} \textsuperscript{(25)} tentesi altında
oturup, perşembe ve pazar günleri çimenler ile tarhların arasına yapılmış
\VUgras{bando yerinde} \textsuperscript{(27)} çalan askerî bandoyu dinlemekten
hoşlanırlar.

%%K t11-c1-08-1 p
8. --- O çimenlerden birinin üstünde tunçtan bir \VUgras{heykel}
\textsuperscript{(28)} yükselir; yurduna büyük hizmetler etmiş bir
şehirlimizin anısına dikilmiş bir anıttır.

%%K t11-tit-3 sub
\VUsaut{4.0mm}
\VUpk{10.2pt}{40}{Gidiş geliş.}
\VUsaut{3.4mm}

%%K t11-c1-09-1 p
9. --- Şehrimiz daha elektrikle aydınlatılmıyor; yalnız \VUgras{havagazı
lambaları} \textsuperscript{(29)} vardır. Ama \VUgras{yerel gazeteler} o
aydınlatma düzeninin iyileştirilmesi için kampanya yürütüyor, tıpkı
tramvayların \VUgras{çekme düzeninin} çoktan iyileştirildiği gibi.
\VUgras{Atların} \VUgras{çektiği} eski arabaların yerine artık kullanışlı
\VUgras{elektrikli tramvaylar} \textsuperscript{(31)} gelmiştir, ki pek
\VUgras{ucuz} bir ücretle yolcuları \VUgras{şehrin} bir ucundan
\VUgras{ötekine} tez ulaştırır.

%%K t11-c1-10-1 p
10. --- Adliyenin yanında \VUgras{banka} \textsuperscript{(32)} vardır, ki
orada ticaret senetleri kırdırılır. \VUgras{Banka tahsildarları}
\textsuperscript{(33)} insanların evine gidip alacakları toplar.

%%K t11-tit-4 sub
\VUsaut{4.0mm}
\VUpk{10.2pt}{40}{Sanat ve bilim.}
\VUsaut{3.4mm}

%%K t11-c1-11-1 p
11. --- Yanındaki o güzel yapı \VUgras{müzedir}. İçinde bir arada şehrin
\VUgras{kitaplığı} \textsuperscript{(34)}, birkaç güzel \VUgras{doğa tarihi
koleksiyonu} \textsuperscript{(35)} (bir doğa tarihi müzesi) ve şık bir
\VUgras{resim galerisi} \textsuperscript{(36)} vardır, ki görkemli bir sürekli
\VUgras{sergi} yapar.

%%K t11-c1-11-2 p
Haftada iki kez halka açılır, ama görmek isteyenler onu herhangi bir gün
gezebilir, \VUgras{bekçiye} \textsuperscript{(37)} küçük bir bahşiş vererek.
\VUgras{Ressamlar} \textsuperscript{(38)} orada çalışmaya gelir. Onları
sehpalarının \VUgras{önünde} görmek olur, elde \VUgras{palet}, ünlü tabloları
kopya ederken.

%%K t11-c1-12-1 p
12. --- Müzenin arkasındaki yükseklerde o \VUgras{kubbe}
\textsuperscript{(40)}, ki \VUgras{hastanenindir} \textsuperscript{(39)}, ile
o
\VUgras{öğretmen okulunun} \textsuperscript{(41)} yapıları güçlükle görünür;
şehir okullarımızın öğretmenleri orada yetişmişti. Tiyatronun alanında bir
\VUgras{postane} \textsuperscript{(43)} vardır, ki bir \VUgras{özel eve}
\textsuperscript{(42)} yerleşmiştir.

%%K t11-tit-5 sub
\VUsaut{5.0mm}
\VUpk{11.4pt}{40}{Yangın.}
\VUsaut{3.0mm}

%%K t11-tit-6 sub
\VUpk{10.2pt}{40}{Yangın çağrısı.}
\VUsaut{3.4mm}

%%K t11-c2-01-1 p
1. --- Şehirde korkunç haber yayılıyor: tiyatroda daha yeni yangın çıkmış!
Ben \VUgras{alana} \textsuperscript{(96)} varana dek kalabalık
\VUgras{kaldırımda} \textsuperscript{(46)} ve \VUgras{caddede}
\textsuperscript{(47)} toplanmıştır. \VUgras{Posta memurları}
\textsuperscript{(44)} ile \VUgras{postacılar} \textsuperscript{(45)}
postaneden çıkıyor. Bir \VUgras{tramvay sürücüsü} \textsuperscript{(48)}
tramvayını durdurmak zorunda kalmıştır, şehrin bir \VUgras{cankurtaranı}
\textsuperscript{(50)} geçebilsin diye; o da bir \VUgras{cerrah}
\textsuperscript{(52)} ile bir \VUgras{hastabakıcı} \textsuperscript{(51)}
getiriyor, bir \VUgras{yaralıya} \textsuperscript{(53)} bakmak için, ki
\VUgras{sedyeyle} \textsuperscript{(54)} \VUgras{gazete kulübesinin}
\textsuperscript{(55)} yanına getirilmiştir.

%%K t11-c2-02-1 p
2. --- Talimden dönen bir piyade bölüğü durmuştur, at üstündeki
\VUgras{şehir muhafızlarıyla} \textsuperscript{(61)} birlikte düzeni korumaya
yardım etsin diye. \VUgras{Atlıların} atları iki ayak üstünde şaha kalkıyor;
\VUgras{askerlerden} \textsuperscript{(57)} ya da \VUgras{jandarmalardan}
\textsuperscript{(63)} daha iyi \VUgras{seyircileri} \textsuperscript{(56)}
geri itebiliyorlar. Jandarmalar zaten pek meşguldür. İçlerinden biri
\VUgras{polis memurlarına} \textsuperscript{(64)} bir \VUgras{yaralıyı}
\textsuperscript{(65)} daha nereye götüreceklerini söylüyor; merdivenden
düşmüş yiğit bir itfaiyeci.

%%K t11-c2-03-1 p
3. --- \VUgras{Subay} \textsuperscript{(66)} tam gelmekte olan
\VUgras{buharlı tulumbanın} \textsuperscript{(67)} nereye konacağını tam tamına
söylüyor. O güçlü aygıtı beklerken bir \VUgras{el tulumbası}
\textsuperscript{(68)} kurdurmuştur, ki uzun \VUgras{hortumu}
\textsuperscript{(69)} ateşin üstüne su seli döküyor. Altı itfaiyeci onu
çalıştırıyor, arkadaşları \VUgras{ağzı} \textsuperscript{(70)} tutup
\VUgras{fışkırtmayı} \textsuperscript{(71)} ateşin ortasına doğrultuyor.

%%K t11-c2-04-1 p
4. --- Tiyatronun öte yanında büyük \VUgras{merdiven} \textsuperscript{(72)}
kurulmuştur. İtfaiyecilerin, kara \VUgras{duman} \textsuperscript{(75)}
burgaçları arasından yükselen o alevlere yaklaşmasını sağlar.

%%K t11-tit-7 sub
\VUsaut{5.0mm}
\VUpk{10.2pt}{40}{Yiğitlik işleri.}
\VUsaut{3.4mm}

%%K t11-c2-05-1 p
5. --- \VUgras{Yangın} \textsuperscript{(74)} \VUgras{tiyatronun}
\textsuperscript{(73)} giriş salonunda çıktı, ilk \VUgras{gösterimi} yarın
olacak o \VUgras{operanın} provası sırasında. Bereket salonda yalnız birkaç
\VUgras{seyirci} \textsuperscript{(76)} vardı. Zavallılar \VUgras{balkona}
\textsuperscript{(77)} sığındı; orada yiğit \VUgras{itfaiyeciler}
\textsuperscript{(86)} merdiven ve halatlarla onlara yardıma geliyor.

%%K t11-c2-06-1 p
6. --- \VUgras{Sundurmanın} \textsuperscript{(78)} altında, ki orada
\VUgras{oyun ilanı} \textsuperscript{(79)} gösterimi duyurur,
\VUgras{çalgıcılar} \textsuperscript{(81)} kaçarken görünüyor, ki
\VUgras{orkestranındır} \textsuperscript{(80)}, çalgılarını taşıyarak:
\VUgras{keman} \textsuperscript{(81)}, \VUgras{flüt} \textsuperscript{(82)},
\VUgras{viyolonsel} \textsuperscript{(83)}, \VUgras{trombon}
\textsuperscript{(84)}, \VUgras{davul} \textsuperscript{(85)}, ve daha
başkaları. \VUgras{Dekorun} \textsuperscript{(87)} bir bölümü kurtarılmaya
çalışılıyor.

%%K t11-c2-07-1 p
7. --- \VUgras{Oyuncular} \textsuperscript{(89)} ile \VUgras{oyuncu kadınlar}
\textsuperscript{(88)}, şarkıcılar ve şarkıcı kadınlar, sahne
\VUgras{giysileriyle} \textsuperscript{(90)} kaçmak zorunda kalmıştır. Baş
şarkıcı bir \VUgras{gazeteciye} \textsuperscript{(92)} bunun nasıl
olduğunu anlatıyor. \VUgras{Muhabir} daha ilk andan yetkililerle birlikte
koşup gelmişti: \VUgras{vali} \textsuperscript{(91)}, \VUgras{belediye
başkanı} \textsuperscript{(93)}, \VUgras{emniyet müdürü}
\textsuperscript{(94)} ve bir \VUgras{yargıç} \textsuperscript{(95)}, ki
sorumluluğun kimde olduğunu soruşturacak.
\VUsaut{6.0mm}
\VUfilet{22mm}
